%%
% This is an Overleaf template for presentations
% using the TUM Corporate Desing https://www.tum.de/cd
%
% For further details on how to use the template, take a look at our
% GitLab repository and browse through our test documents
% https://gitlab.lrz.de/latex4ei/tum-templates.
%
% The tumbeamer class is based on the beamer class.
% If you need further customization please consult the beamer class guide
% https://ctan.org/pkg/beamer.
% Additional class options are passed down to the base class.
%
% If you encounter any bugs or undesired behaviour, please raise an issue
% in our GitLab repository
% https://gitlab.lrz.de/latex4ei/tum-templates/issues
% and provide a description and minimal working example of your problem.
%%

%% ================================== base ================================== %%
\documentclass[
	german,            % define the document language (english, german)
	aspectratio=169,    % define the aspect ratio (169, 43)
	% handout=2on1,       % create handout with multiple slides (2on1, 4on1)
	% partpage=false,     % insert page at beginning of parts (true, false)
	% sectionpage=true,   % insert page at beginning of sections (true, false)
]{tumbeamer}

\usepackage{booktabs}
\usepackage{float}

%% ================================ language ================================ %%
\usepackage{babel}[ngerman]
\usepackage{csquotes}[autostyle]
\usepackage[useregional]{datetime2}
\usepackage{pmboxdraw}

%% ================================ graphics ================================ %%
\usepackage{graphicx}
\usepackage{url}
\usepackage{hyperref}

\graphicspath{ {./resources/} }

%% ================================== tikz ================================== %%
\usepackage{tikz}
\usepackage{pgfplots}

%% ================================ metadata ================================ %%
\author{Niklas Ladurner}
\institute{\theChairName\\\theDepartmentName\\\theUniversityName}

\footline{\insertauthor~|~\insertshorttitle~|~\insertshortdate}

\TUMbeamersetup{
	title page = TUM tower,         % style of the title page
	part page = TUM toc,            % style of part pages
	section page = TUM toc,         % style of section pages
	content page = TUM more space,  % style of normal content pages
	tower scale = 1.0,              % scaling factor of TUM tower (if used)
	headline = TUM threeliner,      % which variation of headline to use
	footline = TUM default,         % which variation of footline to use
	% configure on which pages headlines and footlines should be printed
	headline on = {title page},
	footline on = {every page, title page=false},
}

% available frame styles for title page, part page, and section page:
% TUM default, TUM tower, TUM centered,
% TUM blue default, TUM blue tower, TUM blue centered,
% TUM shaded default, TUM shaded tower, TUM shaded centered,
% TUM flags
%
% additional frame styles for part page and section page:
% TUM toc
%
% available frame styles for content pages:
% TUM default, TUM more space
%
% available headline options:
% TUM empty, TUM oneliner, TUM twoliner, TUM threeliner, TUM logothreeliner
%
% available footline options:
% TUM empty, TUM default, TUM infoline


\title{Übung 07: Profiling und Benchmarking}
\subtitle{Grundlagenpraktikum Rechnerarchitektur}
\date{\DTMdisplaydate{2024}{5}{31}{-1}}

\begin{document}

\maketitle

\begin{frame}[c]{}{}
	\begin{center}
		\LARGE  Keine Garantie für die Richtigkeit der Tutorfolien: Bei Unklarheiten/Unstimmigkeiten
		haben VL-Folien Recht!
	\end{center}
\end{frame}

\begin{frame}[c, fragile]{Profiling und Zeitmessung}{}
	\begin{itemize}
		\item tracing profiler: Eingriffe des Profilers in Programmcode (präzise)
		\item sampling profiler: regelmäßige Samples während Ausführung (ungenauer), bspw. \verb|perf|
		\item Zeitmessung über \verb|clock_gettime(...)|
		\item \verb|CLOCK_MONOTONIC| verwenden!
	\end{itemize}
\end{frame}


\begin{frame}[c, fragile]{Anwendung im Projekt}{}
	\begin{itemize}
		\item \textit{ein} System zum Vergleich von Ausführungszeiten verwenden
		\item am besten bare-metal Linux (keine VMs, keine Rechnerhalle), \verb|-O2|
		\item mehrmals mindestens 1s lang testen, evtl. Programm wiederholen $\rightarrow$ arith. Mittel über Ausführungszeiten
		\item falls Ausarbeitung erforderlich, relevante Systeminformationen angeben (\verb|sudo lshw|)
		\item Begründung, welche Speedups auf welche Optimierungen zurückzuführen sind
	\end{itemize}
\end{frame}

\begin{frame}[c]{}{}
	\begin{center}
		\LARGE Fragen?
	\end{center}
\end{frame}

\begin{frame}[fragile, c]{Links}{}
	\begin{itemize}
		\item Zulip: \href{https://zulip.in.tum.de/#narrow/stream/2267-GRA-Tutorium---Gruppe-20}{\enquote{GRA Tutorium - Gruppe 20}}
		      bzw. \href{https://zulip.in.tum.de/#narrow/stream/2269-GRA-Tutorium---Gruppe-22}{\enquote{GRA Tutorium - Gruppe 22}}
		\item \href{https://gra.caps.in.tum.de}{Praktikumswebsite}
		\item \href{https://www.felixcloutier.com/x86/}{x86 instruction reference by Félix Cloutier}
		\item \href{https://perf.wiki.kernel.org/}{perf Wiki}
	\end{itemize}
\end{frame}

\maketitle

\end{document}
